% Transcription — fonds Alexandre Grothendieck, shelfmark 31, batch 3 (pages 41-60).
% Established from the facsimile archives/batches/31/batch-03.pdf (a mirror of
% https://grothendieck.umontpellier.fr/31.pdf), per the skill
% transcribe-grothendieck.
%
% Pass: Opus 5 (claude-opus-5), 2026-09-11 — first pass, unchecked against
% the pages by a human.
%
% The batch holds three runs. The third is opened by a cover sheet inscribed
% in his hand (page 53), which becomes the section heading and carries no
% \page{}; page 60 is a loose scrap on a different subject and is given its
% own heading.
%
%   Pages 41-44 — ink, fair on 41-43, a fast hand on 44: a proposition on a
%     ring B with a finite group G of automorphisms, A = B^G, the trace
%     Tr_G b = Σ_{g∈G} gb, and the equivalence of a) A = Tr_G B, a') 1 ∈
%     Tr_G B, b) Ĥ(G,M) = 0 for every (G,B)-module M, b(q) the same in one
%     degree; then the relative form for G' ⊂ G with A' = B^{G'} and the
%     transfer t(G,G'), the reduction b'(q) ⟺ b'(0), and the statement that
%     when 1 = Tr_{G/G'} φ for some φ ∈ Hom_{G'}(M,M) the transfer is onto
%     and the corestriction injective with image the stable elements
%     Ĥ(G',M)^G — « Ceci contient le Th. sur les groupes de Sylow ». Page 43
%     applies this to a local ring O with decomposition group G_d and
%     inertia G_i, and closes with the comparison square between Ĥ(G,M),
%     Ĥ(G_i,M)^{G/G_i}, Ĥ(G_d,M)^{G/G_d} and Ĥ(G_d, M_{m_0}). Page 44,
%     cancelled by two long diagonals, is class field theory: the two exact
%     sequences for G_k and G_k^0, their S-versions, and the lemma that an
%     extension of a local field is unramified iff N_{L/k} U_L = U_k.
%   Pages 46-51 — ink, a fast draft: a Proposition giving four equivalent
%     conditions on a field k and an integer n for the vanishing of H^i(k,G)
%     above n — for commutative algebraic groups, for finite groups, and for
%     the groups of units of separable commutative algebras — with the
%     implication diagram and the dévissages (unipotent/semi-simple, the
%     ℓ-adic filtration, characteristic 0 against characteristic p, the
%     passage to the torsion subgroup U of G and H^i(g,U) = lim H^i(g, _nU)).
%     Page 49 proves c) ⟹ b') through A = k̄(Γ), the norm N_{A/k} and the
%     subgroup G_1 of elements of norm 1. Pages 50-51 are two lemmas: a field
%     of transcendence degree 1 over an algebraically closed field is
%     quasi-algebraically closed, with the corollaries for Ĥ^i(K,G_m) and
%     H^i(K,Γ) and the Artin–Schreier sequence for the p-part; then Lemme 2,
%     raising n by one on passing to an extension of transcendence degree 1
%     through the spectral sequence of e → G' → G → G'' → e, and the
%     corollary H^i(K,Γ) = 0 for i > n for a function field of dimension n.
%   Pages 54-59 — ink (59 in pencil), a fast draft under the cover sheet of
%     page 53: does an injective sheaf restrict to an acyclic sheaf on a
%     closed subspace? « Réponse négative. » Four examples build the
%     counter-examples — an irreducible space with two closed points, a
%     circle meeting a line, two lines crossing at a point, two irreducible
%     curves meeting in two points a and b with the Mayer–Vietoris sequence
%     H^0(Y_1) × H^0(Y_2) → H^0(Y_{12}) → H^1(Y,F_Y) → 0 and the section
%     (0,1) ∈ F_a × F_b that does not come from a global one — and page 58
%     lifts the last to a Y of dimension 2 by cutting X with a plane. Page 59
%     is a Proposition listing (i), (i bis), (ii), (ii bis), (iii) on the
%     vanishing of H^i(U,F) and of the Čech groups of a covering, with the
%     implications and the Leray spectral sequence.
%   Page 60 — pencil, a scrap: an N.B. on k → K, the condition k ⊂ ℓ(K^p)
%     for δ(L/K, k/ℓ) = 0, the example L = K(x) with x^p ∈ K and x^p ∉
%     ℓ(K^p), the map Ω^1_{K/k} ⊗_K L → Ω^1_{L/k}, and a boxed inclusion
%     K''_∞ ⊂ K_∞(K_0'^p).
%
% Page 45 is blank and page 52 an empty folder cover; both are absent.
%
% The dating is the inventory's own, for the whole shelfmark.

\documentclass[11pt,a4paper]{article}
% Le préambule commun à toutes les transcriptions du fonds Grothendieck.
%
% Il tient en une idée : **l'incertitude se compose, elle ne se résout pas.**
% Une transcription qui lisse un mot douteux a détruit la seule chose pour
% laquelle on la faisait — un lecteur doit pouvoir savoir, à chaque endroit,
% ce qui était sur le feuillet et ce qui a été ajouté. Les six macros qui
% suivent sont donc l'appareil critique, et non de la décoration.
%
% Les mêmes commandes sont reconnues par `scripts/render.mjs`, qui produit la
% vue de lecture du site. Une macro ajoutée ici doit l'être aussi là-bas,
% faute de quoi le rendu échouera bruyamment — ce qui est voulu : un rendu
% silencieusement faux est pire qu'un rendu absent.

\ProvidesPackage{grothendieck}[2026/08/08 Transcriptions du fonds Grothendieck]

\usepackage{amsmath,amssymb,amsthm}
\usepackage{mathrsfs}
\usepackage{stmaryrd}
% Les diagrammes commutatifs sont partout dans ce fonds : c'est la langue
% même de la géométrie algébrique de Grothendieck, et une transcription qui
% les rendrait en prose ne transcrirait rien.
\usepackage{tikz-cd}
% Français par défaut — c'est la langue des feuillets — mais l'anglais est
% chargé aussi : la traduction et le résumé sont des documents anglais, et la
% typographie française (espaces avant les deux-points) y serait une faute.
% Ces documents basculent par \selectlanguage{english} en tête de corps.
\usepackage[english,french]{babel}
\usepackage{fontspec}
\usepackage{geometry}
\geometry{a4paper,margin=25mm,marginparwidth=28mm}
\usepackage{marginnote}
\usepackage{xcolor}
% ulem, not soul. `soul`'s \st and \ul are built on a character-by-character
% scan that XeLaTeX's font machinery breaks: the document fails with an
% undefined control sequence rather than degrading. ulem does the same two jobs
% and is Unicode-engine safe. `normalem` stops it hijacking \emph, which the
% transcriptions use for Grothendieck's own emphasis.
\usepackage[normalem]{ulem}
\usepackage{hyperref}
\hypersetup{colorlinks=true,linkcolor=black,urlcolor=black,pdfborder={0 0 0}}

% --- Métadonnées du lot -----------------------------------------------------
% Reprises telles quelles par le script de rendu, qui les affiche en tête de
% la vue de lecture. Elles fixent ce que le fichier prétend couvrir : sans
% elles, une transcription flotte sans point d'attache dans un fonds de seize
% mille feuillets.

\newcommand{\folder}[1]{\gdef\grothFolder{#1}}
\newcommand{\batch}[1]{\gdef\grothBatch{#1}}
\newcommand{\leaves}[2]{\gdef\grothFirst{#1}\gdef\grothLast{#2}}
\newcommand{\foldertitle}[1]{\gdef\grothFoldertitle{#1}}
\gdef\grothFolder{?}\gdef\grothBatch{?}\gdef\grothFirst{?}\gdef\grothLast{?}\gdef\grothFoldertitle{}

% --- Le résumé --------------------------------------------------------------
%
% Il ouvre la lecture modernisée, sous le titre « Résumé », et il est composé
% en petit. Ce n'est pas un ornement : c'est le seul passage du document qui
% ne soit pas des mathématiques, et un lecteur doit pouvoir le reconnaître
% avant de l'avoir lu — puis le sauter s'il connaît déjà le sujet, ou s'y
% arrêter s'il ne le connaît pas. Le filet qui le suit marque la reprise du
% corps.

\newenvironment{resume}
  {\small\setlength{\parskip}{0.45em}}
  {\par\medskip\hrule\medskip}

% --- Mots-clés ---------------------------------------------------------------
%
% En anglais, en clôture du résumé de la lecture modernisée : le vocabulaire
% moderne sous lequel chercher ce que les feuillets construisent. Le manifeste
% du site (`npm run manifest`) les extrait pour étiqueter la cote entière —
% cette ligne est la source unique des tags, il n'y a pas de fichier à côté.
\newcommand{\keywords}[1]{\par\smallskip\noindent
  {\footnotesize\textsc{Keywords} --- \textit{#1}}\par}

% --- L'appareil critique ----------------------------------------------------

% Le numéro de feuillet, en marge. C'est l'ancre qui permet de régler un
% désaccord sur une formule en regardant *un* feuillet plutôt qu'une chemise
% de six cents. Le site s'en sert aussi pour tourner les pages du fac-similé
% quand on fait défiler la transcription.
\newcommand{\leaf}[1]{%
  \marginnote{\footnotesize\color{gray!70}\textsf{#1}}%
  \label{leaf:#1}\ignorespaces}

% Illisible : on ne devine pas. Un mot inventé qui se lit comme les autres est
% le pire résultat possible.
\newcommand{\ill}{\textcolor{red!60!black}{[\,\ldots\,]}}

% Lecture douteuse : le mot est donné, et signalé comme incertain.
\newcommand{\uncertain}[1]{\uline{#1}}

% Ajout de l'éditeur — une lettre manquante, un mot sous-entendu.
\newcommand{\add}[1]{\textcolor{blue!50!black}{[#1]}}

% Biffé par Grothendieck lui-même. Ce qu'il a rayé fait partie du document :
% une rature dit souvent où la pensée a changé de direction.
\newcommand{\struck}[1]{\sout{#1}}

% Note du transcripteur, sur l'état matériel du feuillet.
\newcommand{\note}[1]{\par\smallskip{\small\color{gray!60!black}\textit{[#1]}}\par\smallskip}

% Note marginale de Grothendieck, distincte du corps du texte.
\newcommand{\marginal}[1]{\par\smallskip\noindent\hspace{1em}%
  {\small\color{gray!60!black}#1}\par\smallskip}

% --- Titre ------------------------------------------------------------------

\AtBeginDocument{%
  \begin{center}
    {\large\bfseries Fonds Alexandre Grothendieck --- cote n\textsuperscript{o}~\grothFolder}\\[2pt]
    {\itshape\grothFoldertitle}\\[4pt]
    {\small feuillets \grothFirst--\grothLast\ (lot \grothBatch)}\\[2pt]
    {\footnotesize\color{gray!60!black}Transcription établie d'après le fac-similé numérisé
     par l'Université de Montpellier.\\
     Les lectures incertaines sont soulignées, les illisibles notées [\,\ldots\,],
     les ajouts entre crochets.}
  \end{center}
  \vspace{1em}}

\endinput


\folder{31}
\batch{3}
\pages{41}{60}
\dating{[1963-1973]}
\watermark{Édition de démonstration}
\foldertitle{SGA A [SGA 4] (Résidus de rédaction) : notes manuscrites (s.d.)}

\begin{document}

\section*{Groupes finis d'automorphismes, trace et cohomologie de Tate (pages 41 à 44)}

\page{41}
Soit $B$ un anneau avec unité et un groupe fini $G$ d'automorphismes,
$A = B^{G}$. On pose $\operatorname{Tr}_{G} b = \sum_{g \in G} gb$.
Conditions équivalentes \note{Une parenthèse de deux lignes suit l'énoncé,
qui définit $\varphi$ comme l'action de $B$ sur un module et
$\operatorname{Tr}_{G} : B \to A$ ; elle est en grande partie
\ill{}.}
\begin{itemize}
\item[a)] $A = \operatorname{Tr}_{G} B$
\item[a$'$)] $1 \in \operatorname{Tr}_{G} B$
\item[b)] pour tout $(G,B)$-module $M$, $\hat{H}(G,M) = 0$
\item[b(q))] pour tout $(G,B)$-module $M$, $\hat{H}^{q}(G,M) = 0$
\end{itemize}

\textbf{Démonstration} : a) $\Rightarrow$ a$'$) trivial. a$'$) $\Rightarrow$
b) : \ill{} on désigne par $\varphi_{M}(b)$ l'homothétie dans $M$ définie
par $b \in B$, on a $g\,\varphi_{M}(b)\,g^{-1} = \varphi_{M}(gb)$, d'où
\[
\varphi_{M}(\operatorname{Tr}_{G} b) = \operatorname{Tr}_{G}\varphi_{M}(b),
\]
d'où le

\textbf{Corollaire} : Si $\lambda \in \operatorname{Tr}_{G} B$,
l'homothétie $\lambda$ dans $\hat{H}(G,M)$ est nulle
\note{Une incise suit, largement \ill{}, qui observe que la condition est
aussi nécessaire, en prenant $M = B$ ; c'est elle qui donne
a$'$) $\Rightarrow$ b).}

b) $\Rightarrow$ b(q) trivial. On a b(q) $\Rightarrow$ b(q$+1$) car on a la
suite exacte
\[
0 \to M \to M \otimes_{\mathbb{Z}} \mathbb{Z}(G) \to N \to 0,
\qquad m \mapsto \sum_{g} m \otimes g,
\]
où la structure de $B$-module est définie par $\lambda(m \otimes g) =
\lambda m \otimes g$, et la structure de $G$-module par $g'(m \otimes g) =
g'm \otimes g'g$ ; alors $g'(\lambda[m \otimes g]) = (g'\lambda)(g'[m
\otimes g])$ car $g'(\lambda m \otimes g) = g'(\lambda m) \otimes g'g$ et
$(g'\lambda)(g'm \otimes g'g) = (g'\lambda)(g'm) \otimes g'g$,
c'est-à-dire $g'(\lambda m) = (g'\lambda)(g'm)$. \ill{} puisque
$\hat{H}(G, M \otimes_{\mathbb{Z}} \mathbb{Z}(G)) = 0$, on a
$\hat{H}^{q+1}(G,M) = \hat{H}^{q}(G,N) = 0$. \ill{}

\textbf{Corollaire} Supposons $B$ \struck{intègre et intégralement clos},
\marginal{$A$ local.} les conditions de la proposition sont vérifiées si
l'indice de ramification $e$ de $B$ sur $A$ ($=$ ordre du groupe d'inertie)
\struck{\ill{}} est inversible dans $A$, i.e. \ill{}
\struck{$A/\mathfrak{m}(A)$}
\note{Les trois dernières lignes sont biffées et reprises ; seul « est
inversible dans $A$ » se lit avec certitude.}

\marginal{\ill{} le corollaire \ill{} --- Comment voir ça ?}

\page{42}
Avec les conditions et la proposition générale, soit $G' \subset G$, et soit
$A' = B^{G'}$, $A \subset A' \subset B$. Conditions équivalentes
\begin{itemize}
\item[a)] $\operatorname{Tr}_{G/G'} A' = A$
\item[a$'$)] $1 \in \operatorname{Tr}_{G/G'} A'$
\item[b)] pour tout $(G,B)$-module $M$, $t(G,G') : \hat{H}(G',M) \to
\hat{H}(G,M)$ est \struck{injectif} \struck{bijectif} surjectif
\item[b$'$)] \ill{}
\item[b(q))] \ill{}
\item[b$'$(q))] \ill{}
\end{itemize}
\note{Les quatre dernières conditions ne sont pas rédigées : b$'$), b(q) et
b$'$(q) ne sont que nommées, les énoncés restant à écrire.}

a) $\Leftrightarrow$ a$'$) trivial. a$'$) $\Rightarrow$ b) car on a
$t(G,G')\varphi(b) = \varphi(t(G,G')b)$ pour $b \in A'$, d'où si
$t(G,G')b = 1$, $t(G,G')\varphi(b) = 1$, d'où b) d'après un lemme général.

\begin{tikzcd}
\text{b} \arrow[r, Rightarrow] \arrow[d, Rightarrow] & \text{b}' \arrow[d, Rightarrow] \\
\text{b}(q) \arrow[r, Rightarrow] & \text{b}'(q)
\end{tikzcd}

De plus on voit de la même façon que b$'$(q) $\Rightarrow$ b$'$(q$+1$) et
b$'$(q) $\Rightarrow$ b$'$(q$-1$), donc b$'$(q) $\Leftrightarrow$ b$'$(0).
Appliquant b$'$(0) à $M = B$, on trouve que
\[
A'/\operatorname{Tr}_{G'} B \longrightarrow A/\operatorname{Tr}_{G} B
\]
est surjectif, i.e. que l'on a
$1 \in \operatorname{Tr}_{G} B + \operatorname{Tr}_{G/G'} A' =
\operatorname{Tr}_{G/G'} A'$.
\note{La flèche est écrite de droite à gauche sur la page ; elle est ici
rétablie dans le sens de la lecture.}

Reste à prouver : Soient $G' \subset G$ ($G$ fini), soit $M$ un $G$-module
tel que $1 \in \operatorname{Hom}_{\mathbb{Z}}(M,M)$ soit de la forme
$1 = \operatorname{Tr}_{G/G'}\varphi$ \struck{\ill{}} ($\varphi \in
\operatorname{Hom}_{G'}(M,M)$) ; alors
\[
\operatorname{Tr}_{G/G'} : \hat{H}(G',M) \to \hat{H}(G,M)
\ \text{est surjective},
\]
\[
\operatorname{Tr}_{G'/G} : \hat{H}(G,M) \to \hat{H}(G',M)
\ \text{est \uncertain{injective}},
\]
d'image \ill{} l'ensemble des éléments de $\hat{H}(G',M)$ qui sont
« stables », c'est-à-dire $\hat{H}(G',M)^{G}$.
\note{La seconde flèche est également écrite de droite à gauche.}

(Ceci contient le Th.\ sur les groupes de Sylow, \struck{\ill{}} \ill{} en
se plaçant dans la catégorie des groupes \ill{}, mais \ill{} les groupes de
torsion, où par ex.\ \ill{} \ill{}
$\operatorname{Tr}_{G'/G' \cap {}^{g}G'} \ill{}
\operatorname{Tr}_{{}^{g}G'/G' \cap {}^{g}G'}$ \ill{})

\marginal{On peut l'\ill{} en prenant $G' = G_{i}$, mais aussi pour $G'$ un
sous-groupe de Sylow $H$ pour $p$, où $p$ est la caractéristique de $k$. On
a \ill{} : itération. On \ill{} $(\mathcal{O}'^{H}/\operatorname{Tr}_{H}
\mathcal{O}')^{G}$, donc d'abord $M = B$, \ill{} : détermination de
$\mathcal{O}'^{H}/\operatorname{Tr}_{H}\mathcal{O}'$ pour une
\struck{\ill{}} $p$-extension.}

\page{43}
1) Soit $G$ un groupe fini d'automorphismes d'un anneau \uncertain{local}
$\mathcal{O}$. $G$ est transitif sur les idéaux maximaux $\mathfrak{m}$,
soit $G_{d}$ le stabilisateur de $\mathfrak{m}_{0}$, et $G_{i} \subset
G_{d}$ le noyau \ill{} de $G_{d}$, formé des automorphismes qui induisent
l'identité sur $\mathcal{O}/\mathfrak{m}$.

Soit, en conditions
\begin{itemize}
\item[a)] Il existe $x \in \mathcal{O}^{G_{i}}$ tel que
$\operatorname{T}_{G/G_{i}} x = 1$. Donc pour tout $G$-$\mathcal{O}$-module
$M$, on a
\[
\hat{H}(G,M) = \hat{H}(G_{i},M)^{G/G_{i}} .
\]
Donc pour qu'on ait $1 \in \operatorname{T}_{G}\mathcal{O}$, i.e.
$\hat{H}(G,M) = 0$ pour tout $M$, il faut et il suffit que
$[G_{i} : e]$ ne soit pas multiple de la caractéristique $p$ de
$\mathcal{O}/\mathfrak{m}$.
\item[b)] Si $G_{i} = \{e\}$, on a $H^{1}(G, \mathcal{O}^{*}) = \{e\}$.
\end{itemize}

Plus généralement, quel que soit $G_{i}$, l'homomorphisme naturel
$H^{1}(G, \mathcal{O}^{*}) \to H^{1}(G_{i},
(\mathcal{O}/\mathfrak{m})^{*})$ est \uncertain{bijectif} (?). Ceci
\ill{} dans le cas de $GL(n,\mathcal{O})$ --- le \ill{}

(Dans a), en tout \ill{}
\[
\overline{\hat{H}(G,M)} = \hat{H}(G,\bar{M})
= \hat{H}\Bigl(G, \bigoplus_{s \in G/G_{d}} s\,\bar{M}_{\mathfrak{m}_{0}}\Bigr)
= \hat{H}(G_{d}, \bar{M}_{\mathfrak{m}_{0}})
= \overline{\hat{H}(G_{d}, M_{\mathfrak{m}_{0}})}\,.
\]
\marginal{\uncertain{justification} pour Zariski ; il y a donc \ill{} hyp.\
\ill{} : finie ??}

Considérons donc le diagramme :

\begin{tikzcd}[column sep=small, row sep=small, nodes={font=\scriptsize}]
\hat{H}(G,M) \arrow[r, "\sim"] \arrow[d, "?"'] &
\hat{H}(G_{i},M)^{G/G_{i}} \arrow[r, "\alpha"] &
\hat{H}(G_{i}, M_{\mathfrak{m}_{0}})^{G_{d}/G_{i}} \\
\hat{H}(G_{d},M)^{G/G_{d}} \arrow[r, "\beta"'] &
\hat{H}(G_{d}, M_{\mathfrak{m}_{0}}) \arrow[r, "\sim"'] &
\hat{H}(G_{d},M)_{\mathcal{O}_{d} - \mathfrak{m} \cap \mathcal{O}_{d}}
\arrow[u, "?"']
\end{tikzcd}

\marginal{$\alpha$ et $\beta$ sont injectifs, à images denses ; on passe aux
\ill{} !}

\marginal{\ill{} de dim ?, $\mathcal{O} \to$ \ill{} $\hat{H}(G,M)$ ?
\ill{}}

\page{44}
\note{Toute la page est barrée de deux longues diagonales qui se croisent.}
\[
0 \to G_{k}^{0} \to G_{k} \to \mathfrak{P}_{k} \to 0
\]
\note{Sous $G_{k}$ : « extension abélienne \ill{} maximale » ; sous
$\mathfrak{P}_{k}$ : « extension non ramifiée maximale ». Le second groupe
est écrit d'un $P$ barré ; la lecture $\mathfrak{P}_{k}$ est conventionnelle.}
\[
0 \to \mathbb{Z}_{2}^{r_{1}} \to G_{k}^{0} \to
\mathcal{U}_{K}^{0}/\operatorname{Im} E \to 0
\]
\note{Sous $\mathcal{U}_{K}^{0}$ : « idèles unités \ill{} » ; sous
$\operatorname{Im} E$ : « idèles unités \ill{} principaux ».}

\textbf{N.B.} $G_{k}/\mathbb{Z}_{2}^{r_{1}}$ \uncertain{pourrait} être un
groupe dont la définition ne fait plus intervenir \ill{} que le groupe de
Galois de l'extension abélienne non ramifiée \ill{} maximale (i.e. qui est
\uncertain{complètement décomposée} au-dessus de toutes les places réelles
de $k$) --- soit de façon générale, $S$ un ensemble \uncertain{fini} de
places, finies ou infinies réelles, alors le groupe de Galois de
l'extension abélienne non ramifiée hors $S$ maximale, soit $G_{k}^{S}$,
\[
0 \to G_{k}^{0,S} \to G_{k}^{S} \to \mathfrak{P}_{k} \to 0
\]
\[
0 \to \mathbb{Z}_{2}^{S_{\infty}} \Big/ \bigcap_{k}
\operatorname{Im}\bigl(E^{(S,k)}\bigr)
\to G_{k}^{0,S} \to \mathcal{U}_{K}^{S,0}/\operatorname{Im} E \to 0
\]
\note{Sous la seconde suite, au crayon : « ceci n'est \ill{} intéressant
$\overline{S^{\infty}}$ ??? $\mathbb{Z}_{2}$ ».}

Ceci résulte du lemme : Pour qu'une extension \struck{abélienne} $L$
\struck{\ill{}} d'un corps local \struck{\ill{}} $K$ soit non ramifiée, il
faut et il suffit que $N_{L/k}\mathcal{U}_{L} = \mathcal{U}_{k}$, ou
$N_{L/k} L^{*} \supset \mathcal{U}_{k}$.

\section*{Une proposition sur $H^{i}(k,G)$ et deux lemmes sur les corps de fonctions (pages 46 à 51)}

\page{46}
\textbf{Proposition} Soit $k$ un corps, $n$ un entier. Conditions
équivalentes :
\begin{itemize}
\item[a)] $H^{i}(k,G) = 0$ pour $i > n$ et tout groupe alg.\ commutatif $G$
défini$/k$
\item[b)] $H^{i}(k,\Gamma) = 0$ \struck{\ill{}} pour $\Gamma$ fini défini$/k$
\item[b$'$)] $H^{i}(k,\Gamma) = 0$ pour $i > n$ et \struck{tout groupe fini
qui est \ill{}} \add{(fini)} \ill{} \struck{\ill{}} sur un corps
$\mathbb{Z}/\ell$, donc le \ill{} $\mathbb{V}/k$ qui est (ou un groupe
abélien) \uncertain{irréductiblement} \ill{}
\item[c)] $H^{i}(k,G) = 0$ pour $i > n$ et tout p.-g.\ alg.\ commutatif $G$
défini$/k$ \struck{\ill{}} devenant isomorphe à un groupe \struck{\ill{}}
sur la clôture galoisienne de $k$, $(\mathbb{G}_{m})^{n}$ \ill{} ---
chaque fois que $G$ est le groupe des éléments inversibles d'une algèbre
\add{(commut.)} séparable \struck{\ill{}} déf$/k$
\item[c$'$)] $H^{i}(k,G) \ne 0$ pour $i > n$, \ill{}
\end{itemize}
\note{L'énoncé est un brouillon rapide : b$'$) et c) sont écrits par-dessus
des versions biffées, et la prose qui les relie est en grande partie
perdue.}

\textbf{Démonstration.}

\begin{tikzcd}[column sep=small, row sep=small, nodes={font=\scriptsize}]
& \text{a} \arrow[dl, Rightarrow] \arrow[dr, Rightarrow] & \\
\text{b} \arrow[d, Rightarrow] & & \text{c} \\
\text{b}' & &
\end{tikzcd}

Prouverons c) $\Rightarrow$ b$'$) et b$'$) $\Rightarrow$ b), ce qui
achèvera la démonstration.

b$'$) $\Rightarrow$ b) : \struck{\ill{}} comme on voit par un
\struck{rapport} \add{dévissage} d'abord de $\Gamma$ \struck{\ill{}}
\add{commutatif}, dévissage immédiat en considérant la filtration
\[
\Gamma \supset \ell\Gamma \supset \ell^{2}\Gamma \supset \cdots
\supset \ell^{k}\Gamma = 0 ,
\]
\ill{} par $\ell$, i.e. \ill{} \struck{\ill{}}

b) $\Rightarrow$ a) : En effet, \struck{Considérons} \ill{} Considérons
\ill{} le plus grand \ill{} groupe unipotent de $G$, qui est défini sur
$k$, \ill{} de sorte que
\[
0 \to G_{u} \to G \to G/G_{u} \to 0 .
\]

\page{47}
Cela nous ramène au cas où $G$ est soit \uncertain{complet}, soit
\uncertain{linéaire}. Dans le deuxième \struck{\ill{}} cas,
\struck{\ill{}} on a $G = G_{u} \times G_{ss}$ (décomposition en partie
semi-simple et unipotente) et cette décomposition est définie sur $k$. Donc
on peut se ramener au cas où $G$ est soit unipotent, soit
\struck{semi-simple} \ill{} d'éléments semi-simples, soit commutatif
\ill{} complet. Dans les deux derniers cas, soit $\mathcal{U}$ \ill{} le
sous-groupe des éléments d'ordre fini de $G = G^{\bar{k}}$ ; alors le
groupe $G/\mathcal{U}$ est un \ill{}

\marginal{ici c'est faux : le \ill{} inductive \ill{} commutatif}

\ill{} $G_{0}$ (comp.\ connexe de $e$) \ill{} variété abélienne, et c'est
un th.\ de Weil pour \ill{} (i.e. $G \simeq \bar{k}^{*n}$) \ill{}
\struck{\ill{}} $G \simeq G_{0} + G/G_{0}$, \ill{} un quotient de
$\mathbb{V}/k$, d'où \ill{} $G/G_{0}$ \ill{}

\ill{} on voit que $G/\mathcal{U}$ s'identifie au quotient de $G_{0}$ par
$\mathcal{U} \cap G_{0}$, donc \ill{} est divisible. Donc
$H^{i}(\mathfrak{g}, G/\mathcal{U}) = 0$, \struck{tandis que} d'où pour
$i > 0$
\[
H^{i}(\mathfrak{g}, G) \simeq H^{i}(\mathfrak{g}, \mathcal{U})
\quad \text{pour } i \geq 1
\]
\note{$\mathfrak{g}$ désigne ici le groupe de Galois ; il est écrit d'un
$g$ à longue hampe, et la page passe de $H^{i}(k,-)$ à
$H^{i}(\mathfrak{g},-)$ sans le dire.}

(comme \ill{} pour $G \to G/\mathcal{U}$ est surjectif \ill{}) \ill{}
$\mathcal{U}$, le \ill{}-groupe de $\mathcal{U}$ formé des éléments
d'ordre premier \ill{}. On a
\[
H^{i}(\mathfrak{g}, \mathcal{U}) = \varinjlim_{n}
H^{i}(\mathfrak{g}, {}_{n}\mathcal{U})
\]

\page{48}
et lorsque $i > n$, les \struck{deux} $H^{i}(\mathfrak{g},
{}_{n}\mathcal{U})$ sont nuls (car ${}_{n}\mathcal{U}$ est fini) ---
d'après b), d'où résulte notre assertion.

\ill{} $G$ un groupe linéaire unipotent commutatif, on distingue deux cas :

1$^{\circ}$) Caractéristique nulle. Alors $G$ \struck{est} $\simeq
\bar{k}^{n}$ sur $\bar{k}$, donc aussi sur $k$ (car $H^{1}(k, Gl(n)) = 0$).
Et \struck{comme} $H^{i}(k, \mathbb{G}_{a}) = 0$ pour $i > 0$, l'assertion
est valable.

2$^{\circ}$) Car.\ $p \ne 0$. Alors on a la filtration
\[
G \supset pG \supset \cdots \supset p^{\ell}G = 0 .
\]
On \ill{} ramener au cas \add{i.e. où} $pG = 0$ (car chaque $p^{i}G$ est
\emph{défini sur $k$}). Mais, on sait \ill{} alors $\mathbb{G}_{a} \simeq
\bar{k}^{n}$ sur $\bar{k}$, donc \ill{} aussi sur $k$ (car
$H^{1}(k, Gl(n)) = 0$), et on \ill{} comme dans 1$^{\circ}$.

Supposons que $k$ soit quelconque \ill{}
\note{La page s'arrête sur cette phrase, laissée en suspens.}

\page{49}
c) $\Rightarrow$ b$'$) : Soit en effet $\Gamma$ un groupe comme dans
b$'$) ; supposons d'abord $\ell \ne p$. \struck{Il suffit} \add{\ill{} dans
la clôture \ill{} de $\Gamma$ \ill{} évidente.} Soit $A = \bar{k}(\Gamma)$
\struck{l'algèbre} \add{une algèbre} \ill{} \add{c'est une algèbre déf$/k$,
commutative} et \struck{semi-simple} séparable (car $[\Gamma : e]$ est
premier à $p$), donc \ill{} on a $H^{i}(k,G) = 0$ si $i > n$, $G$
désignant le groupe des éléments inversibles de $A$. Soit $G_{1}$ le
sous-groupe de $G$ formé des $x \in G$ tels que $N_{A/k}\,x = 1$
(i.e. \struck{soit} $N_{A/k}\,x = $ dét.\ de la translation à gauche par
$x$ dans $A$) ; c'est un sous-groupe défini sur $k$ (noyau de
l'homomorphisme $N_{A/k} : G \to \bar{k}^{*}$, qui est déf$/k$).

Comme $\Gamma \to \ill{}$ \ill{} $\mathbb{V}$-module, \add{et comme
$\Gamma \subset G$} il est visible \add{que} \struck{\ill{}} l'on a
$\Gamma \subset G_{1}$, i.e. $\Gamma \cap G_{1} = \struck{\ill{}}
\{e\}$, et on vérifie facilement que \ill{} (distinguer le cas $\ell$
premier impair, et le cas $\ell = 2$ ; \ill{} dans ce dernier cas \ill{}
et $N_{A/k}\,\xi = \pm 1$ si $\xi \in \Gamma$).

Comme $\ell\Gamma \struck{\ill{}} = \{0\}$, \ill{} enfin $\Gamma \subset
G_{1}$. Il est d'ailleurs visible que $\Gamma \to \ill{} G_{1}$. Le noyau
\ill{} de l'hom.\ $G_{1} \xrightarrow{\ \ell\ } G_{1}$ \ill{} d'ordre
$\ell^{n_{1}}$, i.e. $n_{1} = n - 1$ \ill{}
\note{La fin de la page est une suite de calculs rapides dont la prose ne
se lit pas.}

\page{50}
\textbf{Lemme 1} Soit $k$ un corps algébriquement clos, $K$ une extension
de type fini, de degré de transcendance $1$, de $k$. Alors $K$ est
\struck{\ill{}} quasi-algébriquement clos.

\textbf{Corollaire 1} $\hat{H}^{i}(K, \mathbb{G}_{m}) = 0$ pour tout $i$.

En effet, $K$ est un corps \ill{} galoisienne $K'$, \ill{} quasi-alg.\
clos, donc \struck{\ill{}} le corollaire est bien connu.

\textbf{Corollaire 2} $H^{i}(K,\Gamma) = 0$ \struck{pour $i > 0$} pour
$i > 1$ et \ill{} $\Gamma$ d'ordre fini, i.e. $G_{K}$ opère trivialement.

Soit \add{une} \struck{\ill{}} clôture galoisienne de $K$, $\tilde{K}$.
\struck{\ill{}}

\note{Suit un bloc encadré, écrit très vite, dont voici ce qui se lit :
$\mathcal{U} \subset \tilde{K}^{*}$ le sous-groupe des racines de l'unité
dans $\tilde{K}$ ; alors $\tilde{K}^{*}/\mathcal{U}$ est uniquement
divisible si \ill{} ; d'où \ill{} ; $\hat{H}^{i}(K, \ill{})
\xrightarrow{\sim} \hat{H}^{i}(K, \mathbb{G}_{m})$ \ill{} mod $p$
« isomorphismes ? » ; le deuxième est nul, donc le premier \ill{}.}

Plus généralement \ill{} la caractéristique, alors
$\tilde{K} \xrightarrow{\,x \mapsto x^{n}\,} \tilde{K}$ \ill{} est un
épimorphisme de noyau $\simeq \mathbb{Z}/n$, donc la suite exacte de
cohomologie \ill{} on a donc $H^{i}(K, \mathbb{Z}/n) = 0$. Reste \ill{}
$H^{i}(K, \mathbb{Z}/p) = 0$ si $i > 1$ : \ill{} on \ill{}
\[
0 \to \mathbb{Z}/p \to \tilde{K} \xrightarrow{\ \wp\ } \tilde{K} \to 0 .
\]

\page{51}
\textbf{Lemme 2} Soit $k$ un corps, \add{$\Gamma$ (?)} \struck{\ill{}}
tel que $H^{i}(k',\Gamma) = 0$ pour toute extension $k'$ \add{\ill{}} de
type fini de $k$, $i > n$. \struck{Soit $K$ une extension de degré de
transcendance $\le 1$.} Alors $H^{i}(K,\Gamma) = 0$ pour $i > n+1$.

\begin{tikzcd}[column sep=small, row sep=small, nodes={font=\scriptsize}]
k' = \tilde{k}k & k \arrow[l] \\
K \arrow[u] & \tilde{k}K \arrow[l] \arrow[r, "G'"] & \tilde{K}K
\end{tikzcd}
\note{Le diagramme de la page est un croquis d'extensions ; il est ici
redressé, et les groupes $G$, $G'$, $G''$ y sont portés le long des
flèches auxquelles ils correspondent.}

d'où une suite exacte
\[
e \to G' \to G \to G'' \to e ,
\]
d'où une suite spectrale
\[
H^{i}\bigl(G'', H^{j}(G',\Gamma)\bigr) \Longrightarrow H^{*}(G,\Gamma).
\]

Or $H^{j}(G',\Gamma) = 0$ si $j > 1$ en vertu du lemme 1, et d'autre part
$H^{i}\bigl(G'', H^{j}(G',\Gamma)\bigr) = 0$ si $i > n$, en
\struck{$G''$ étant \ill{}} \ill{} d'après \struck{fini de
l'hypothèse}. D'où \struck{\ill{}} $H^{*}(G,\Gamma) = 0$ si $* > n+1$,
cqfd.

\textbf{Corollaire} Soit $K$ un \struck{\ill{}} corps de fonctions
algébriques sur un corps alg\add{ébriquement} clos $k$,
$\deg\operatorname{tr}_{k} K = n$. Alors $H^{i}(K,\Gamma) = 0$ pour
$i > n$ et \add{d'ordre fini} tel que \ill{} $\Gamma$, i.e. $G_{K}$ opère
trivialement.

\section*{Th.\ fondamentale des morphismes propres [Changement de base] : Théorème d'échange --- Questions marginales (pages 54 à 59)}

\note{Le titre est celui de la page 53, un feuillet de couverture inscrit
de sa main et par ailleurs vierge, qui ne reçoit donc pas de numéro de page
propre. Une ligne biffée le précède, illisible.}

\page{54}
\textbf{Question} $X$ espace topologique, $Y$ fermé de $X$, $F$ faisceau
\struck{injectif} sur $X$, a-t-on
\[
H^{i}(Y, F|Y) = 0 \quad \text{pour } i > 0
\]
\[
\bigl[H^{i}(V, F|V) = 0 \ \text{pour tout \ill{} voisinage } V
\text{ de } Y,\ i > 0\bigr].
\]
Réponse négative.

Tout d'abord, il n'est pas vrai que pour un faisceau $F$ sur $X$ injectif,
l'homomorphisme
\[
\varinjlim_{U \supset Y} H^{0}(U,F) \longrightarrow H^{0}(Y,F)
\]
soit \uncertain{surjectif}, i.e. qu'une section \ill{} se prolonge à un
voisinage [i.e. si $X$ \ill{} $2^{\circ}$ $F$ est \ill{} injectif].

\textbf{Ex 1} $X$ espace irréductible, $Y$ fermé de \ill{} deux points
fermés distincts $a$, $b$ ; $F$ faisceau des \ill{}, la section de $F|Y$
qui vaut $0$ en $a$, la valeur $1$ en $b$, ne se prolonge pas
\struck{à un voisinage} \ill{} $X$, car \ill{} contient un $U$ de $a$
\ill{} un voisinage $V$ de $b$, et $U \cap V \ne \emptyset$, absurde.

\textbf{Ex 2} $X = \overline{M}$ \ill{} une courbe algébrique, $Y$ \ill{}
irréductible, \add{\ill{}} \struck{\ill{}}, mais \ill{} par exemple une
\ill{} $H^{0}$ \ill{} de Zariski.
\note{La fin de l'exemple 2 est un brouillon rapide dont la prose ne porte
plus.}

\page{55}
\note{La page ouvre sur une figure : un cercle, coupé par une droite $Y$ en
deux points $a$ et $b$ ; l'ensemble est $X$.}

Considérons $X_{1} = $ cercle de la figure, et le faisceau $F_{1}$
\add{sur $X$} \struck{\ill{}} induit dans $\mathbb{Q}$, soit $F$ son image
sur $X$. La section \ill{} prolonge une section \ill{} elle \ill{} enfin
sur $Y \cap X_{1}$ à $X_{1}$, c'est en général impossible \ill{}

\textbf{Ex 3} Un faisceau injectif \struck{\ill{}}
\note{L'exemple 3 est barré de plusieurs longs traits horizontaux. On y
distingue les exemples 1 et 2 rappelés, $F|Y$, $X$, $Y$, et un
$\mathbb{Q}$ ; le reste ne se lit pas.}

On ne s'en tire pas \ill{} $Y$ \ill{} un compact irréductible [i.e. \ill{}
le faisceau irréductible n'est \ill{} de type fini \ill{}]. \ill{} $Y$ le
voisinage \ill{} $V$ \ill{}, sur $X$ irréd.

Mais prenons pour $Y$ $Y_{1}$ et $Y_{2}$ deux \ill{} \ill{}

\page{56}
\note{Une figure ouvre la page : deux droites $Y_{1}$ et $Y_{2}$ se
coupant en un point $a$, accolées par une accolade portant $Y$.}

Soit $F$ le faisceau \struck{\ill{}} des \ill{} \ill{} dans $\mathbb{Q}$,
et $Y' = Y - \{a\}$ ; considérons la section de $F|Y'$ qui vaut $0$ sur
$Y'_{1}$, $1$ sur $Y'_{2}$ ; \ill{} \ill{} un $U$ de $a$, et un \ill{}
$\mathbb{Q}$ qui \ill{} voisinage de $Y_{1}$ ($Y_{2}$), absurde car \ill{}
\ill{} !

\note{Suit un bloc entièrement barré, sur lequel se lisent
$Y'_{1}$ ($Y'_{2}$) induisant $0$ (1), puis, plus bas, la suite}
\[
\struck{F_{Y} \to \textstyle\coprod_{i} F_{Y_{i}} \to G \to \ill{}}
\]
\[
\struck{H^{0}(Y_{2}, F_{Y_{i}}) \to H^{1}(Y,G) \to H^{1}(Y, F_{Y}) = 0}
\]
\note{avec, dans la marge du bloc, « remarquer que $H^{1}(Y, F_{Y_{i}}) =
0$ ».}

\page{57}
\textbf{Exemple 4} Donnons un contre-exemple à la seconde question en \ill{} : $Y$ \ill{} \ill{} dans une courbe irréductible,
$Y_{1}$, $Y_{2}$ irréductibles se coupant en deux points $a$, $b$, le $X$
irréductible. $F$ est le faisceau des \ill{} de $X$ dans $\mathbb{Q}$. On a
\[
H^{0}(Y_{1}, F_{1}) \times H^{0}(Y_{2}, F_{Y_{2}}) \to
H^{0}(Y_{12}, F_{12}) \to H^{1}(Y, F_{Y}) \to 0
\]
\[
[\,F_{i} = F_{Y_{i}},\quad F_{12} = F_{Y_{12}}\,]
\]
\note{La page ouvre sur une figure : deux ovales $Y_{1}$ et $Y_{2}$ qui se
recoupent, leurs deux points communs marqués $a$ et $b$.}

\struck{\ill{}} $F_{a} \times F_{b}$ dans le premier \ill{} \ill{} car
\[
H^{0}(Y_{12}, F_{12}) = F_{a} \times F_{b}
\]
\ill{} $f_{1} - f_{2}$, i.e. $f_{i} \in H^{0}(Y_{i}, F_{i})$. Or
$(0,1) \in F_{a} \times F_{b}$ \ill{} \ill{} $\mathcal{U}$.

En effet, si on avait une \ill{} $f_{1}$, $f_{2}$ \ill{}
\note{Suit un bloc barré où l'on distingue $\varphi_{1}, \varphi_{2} \in
\Gamma(U,F)$, $f_{ix} = \varphi_{ix}$ et $f_{ix} = \varphi'_{ix}$, puis
$\varphi'_{2} - \varphi'_{1} = 1$ ; la prose entre les formules est
perdue.}

\page{58}
$x \in U \cap U' \cap Y_{1}$, donc dans un voisinage $W_{1}$ de
$U \cap U' \cap Y_{1}$, et $\varphi_{2x} = \varphi'_{2x}\ (= f_{2x})$ pour
$x$ \ill{} ; \struck{i.e. $\varphi_{1x} = \varphi$} pour $x \in U \cap U'
\cap Y_{2}$, donc
\[
\varphi_{1x} = \varphi'_{1x} + 1 \quad \text{pour } x \in U \cap U'
\cap Y_{2},
\]
\ill{} pour $x$ dans un voisinage $W_{2}$ de $U \cap U' \cap Y_{2}$. Donc
si $x \in W_{1} \cap W_{2}$, \ill{}, et
\[
\varphi_{1x} = \varphi'_{1x} + 1, \qquad \varphi_{1x} = \varphi'_{2x},
\]
i.e. $0 = 1$, absurde [compte tenu \ill{} $W_{1} \cap W_{2} \ne
\emptyset$, $X$ étant irréductible].

Utilisant \ill{} exemple, on trouve un exemple avec $Y$ \emph{irréd.}
(de \emph{dim 2} \ill{}), \ill{} ceci, dans $X$ \ill{} $X_{1}$ plan
(dim 2) qui, considérant un $X_{1}$ plan \ill{} recoupe $Y$ suivant $2$
\ill{}
\[
Y_{1} = X_{1} \cap Y .
\]
\note{Deux figures accompagnent le passage : un parallélogramme portant
$Y$, et deux ovales qui se recoupent.}

On prend alors le faisceau $F$ \ill{} sur $X$ \ill{} des \ill{} de $X_{1}$
dans $\mathbb{Q}$.

\page{59}
\note{Page au crayon, d'une écriture rapide ; les conditions sont
numérotées mais leurs énoncés ne sont écrits qu'à demi.}

\textbf{Proposition} $F$ faisceau abélien sur \ill{} $V$ \ill{} tout $V$
(ordinaire) $X$. Conditions :
\begin{itemize}
\item[(i)] \ill{}
\item[(i bis)] Pour tout ouvert $U$ de $X$, $H^{i}(U,F) = 0$
\item[(ii)] Pour tout ouvert $U$ de $X$, $H^{i}(U,F) = 0$ pour $i > 0$
\item[(iii)] Pour tous ouverts $U$, $V$ de $F$ \add{[pour $X$ noeth.]}
\end{itemize}
\[
H^{1}\bigl(C(U,V); F\bigr) = H^{0}(U \cap V, F) \big/
\operatorname{Im}\bigl(H^{0}(U,F) \times H^{0}(V,F)\bigr)
\]
\begin{itemize}
\item[(ii bis)] Pour tout \uncertain{recouvrement} $U$ et tout
recouvrement $\underline{\mathcal{U}} = (\mathcal{U}_{i})$ de
$\mathcal{U}$, \ill{} $H^{i}(\underline{\mathcal{U}}, F) = 0$ pour $i > 0$.
\end{itemize}

(i bis) $\Rightarrow$ (ii bis) par suite spectrale de Leray, et bien connu
(petit argument de Cartan \ill{} immédiat)

(i bis) $\Rightarrow$ (i) trivial

(ii bis) $\Rightarrow$ (ii) trivial

(i) $\Leftrightarrow$ (ii) car $H^{1}(U \cap V, F) \subset
H^{1}(U \cup V, F)$, d'où \ill{} ; pour (ii) $\Rightarrow$ (i), \ill{} en
considérant $\xi \in H^{1}(U,F)$ \add{(N.B. $X$ \ill{})} \ill{}
$U' \subset U$ \ill{} $\xi$ soit nul, \ill{} dis que $U' = U$, \ill{}

\section*{N.B. : $N_{L/K}$ et la condition $k \subset \ell(K^{p})$ (page 60)}

\page{60}
\note{Feuillet au crayon, sans rapport apparent avec ce qui précède, barré
de deux longs traits obliques.}

\textbf{N.B.} $k \to K$ \ill{}, \ill{} \ill{} que toute extension $L$ de
\ill{} (ou finie) de $K$, \ill{} $N_{L/K}$
\[
\delta(L/K,\ k/\ell) = 0 \quad \text{i.e.}
\]
$k$ soit $\ell$-\ill{} \ill{} $L/K$, il faut et il suffit que
\[
k \subset \ell(K^{p}) \qquad \text{i.e.} \quad k(K^{p}) \subset
\ell(K^{p})
\]
\ill{}
\[
K \to L = K(x), \qquad x^{p} \in K,\quad x^{p} \notin \ell(K^{p})
\]
\[
N_{L/K/k} = 0, \qquad N_{L/K/\ell} \simeq L, \qquad N_{L/K}
\]
\[
\Omega^{1}_{K/k} \otimes_{K} L \longrightarrow \Omega^{1}_{L/k}
\]
\[
K_{0} \ \big| \ K \qquad k(K_{0}^{p}) \subset K_{0}^{p}
\]
\struck{$k \subset$} \struck{$K_{\infty}$ \ill{}} $k$
\[
\boxed{\ K''_{\infty} \subset K_{\infty}(K_{0}'^{p})\ }
\]
\struck{$K_{\infty}(\ill{}) \subset$} \ill{}

\end{document}
